% ============================================================================
% SEC 6 v7 (2026-08-22). Fred's v6 edits, compressed. Two reverts flagged in
% chat: NSR headroom is rejection/p_fail, not bias (5.3; D.5 only if
% truncated). 2.3x removed.
% ============================================================================
\newcommand{\owed}[1]{{\color{red}\textbf{[owed: #1]}}}

\section{Evaluation}
\label{sec:eval}

The emulator closes Sec.~5: what it runs is the compiler's output, so every
number here is a measurement of the estimator the compiler ships, executed
under a device-level error model (Appendix~C.3) and measured against the
gradient of the objective the device realizes, $\nabla C_{\mathrm{device}}$
(Appendix~C). The emulator's cost is exponential in qubits, so execution
results are at few qubits; the compiler's is not, and the device-ready
schedules it emits for every evaluated program, at sizes beyond a hundred
qubits, are shown in Appendix~F. The workload is the running instance of
Sec.~5.2, grown: two qubits where one gradient is examined in detail,
seven where program structure is varied, a thousand where the compiler is
timed. The section follows the order in which the earlier sections made
their promises: the theorems hold on the machine, the fallback does not,
the two sound lowerings are a real choice, and the choice is still there
at device scale.

\subsection{The Estimator on the Machine}
\label{sec:eval-machine}

Theorem~4.3 and its noisy forms (Theorems~C.7 and~C.8) say nothing about a
device; the first check is that they survive one. Figure~\ref{fig:F6} is
that check. Compiled NSR and PSR, each branch executed under the full error
model, converge to $\nabla C_{\mathrm{device}}$ at the $N^{-1/2}$ of
Corollary~5.2 over four decades of budget. Nothing in the lowering ---
normal form, solver, schedule, ledger --- has introduced bias, which is at
once the soundness theorems measured and the correctness evidence for
Sec.~5. The one bias the paper does not remove is on the same axes: PSR
with a real gate floors where Lemma~C.9 puts it, first order in the gate's
infidelity, with nothing dividing by a small step.

\begin{figure}[t]
  \centering
  \includegraphics[width=\linewidth]{figs/F6.pdf}
  \caption{Compiled PSR and NSR converge to $\nabla C_{\mathrm{device}}$ at
  the $N^{-1/2}$ rate over four decades of budget; finite differences floor
  at their deterministic control-precision bias, both at a
  oracle-tuned $\varepsilon^{\ast}$ and at a fixed deployable $\varepsilon$.
  Inset: RMSE against step size
  $\varepsilon$ at fixed budget---the finite-difference ``V'' floors above
  both flat sound-strategy references, and $\times$ mark the $\varepsilon$ at
  which FD's gradient takes the wrong sign (points uphill). TFIM coupling
  gradient, device model of Appendix~C.3, $T/T_2^{\ast}{=}0.15$.}
  \label{fig:F6}
\end{figure}

With the sound estimators established, the same figure settles the
question Sec.~2 opened. Figure~1 showed the finite-difference trap at one
point of one parameter at the Hamiltonian level; the objection is that a
tuned step on the real artifact might escape it. Figure~\ref{fig:F6} gives
FD the step no experimenter has --- chosen by retrospective global sweep
--- and it floors anyway, above both sound strategies, because the
control-precision bias is deterministic at an operating point and
averaging cannot see it (inset: the floor is the minimum over every step
the machine can dial. The instance is chosen
for this: its landscape has features narrower than the dialable steps, as
landscapes do near convergence, where optimization spends its time. On a
smooth landscape FD is fine and competitive at low budget; what it cannot
do is stay fine as the optimizer approaches a minimum.

A gradient-error floor earns its place only if the optimizer notices.
Figure~\ref{fig:loop} runs the two-qubit instance in closed loop, every
method on one optimizer and one budget. NSR and PSR reach the optimum and
hold it; oracle-tuned FD settles a systematic distance short, and a step
one notch larger wanders off. The stall distance is gradient bias over
local curvature --- not the control resolution itself --- which is why
different steps stall in different places and why no step fixed before the
run prevents it. This is
the failure Sec.~2 predicted, now on the shipped artifact, and it is why
the compiler emits Sec.~4's strategies and keeps Sec.~4.4 as a baseline.

\begin{figure}[t]
  \centering
  \includegraphics[width=\linewidth]{figs/F_loop.pdf}
  \caption{The two-qubit instance in closed loop: median\,$\pm$\,IQR of
  $\lVert\theta_t-\theta^{\ast}\rVert$ over 20 seeds, one optimizer and one
  shot budget for every method, on a deliberately ill-conditioned (sloppy)
  self-consistent recovery valley whose interior minimum $\theta^{\ast}$ is
  fixed by construction. The $\varepsilon$-free strategies PSR and NSR reach
  $\theta^{\ast}$ and hold it within tolerance; the finite-difference steps
  fail in three ways---too large diverges (truncation), too small is
  unreliable (wide IQR), and the retrospectively best $\varepsilon$ floors
  just above tolerance. Filled markers: the first step from which the median
  stays inside tolerance for five consecutive steps. Inset: the valley, with
  each method's median path (star $=\theta^{\ast}$, open circle $=$ start;
  the too-large-$\varepsilon$ path runs off the frame).}
  \label{fig:loop}
\end{figure}

\subsection{Pricing the Pair}
\label{sec:eval-select}

Two sound lowerings remain, and Sec.~5.3 priced them without choosing:
NSR pays statistics and headroom, PSR pays motion and a gate, with the
constants of Corollary~5.2. If one always won, the strategy interface of
Sec.~4 would be an abstraction without a use. Figure~\ref{fig:select}, on
the seven-qubit instance, shows it is not: on the plane the compiler sees
--- how many coefficients a program differentiates against how many terms
each touches, both read off program text --- PSR's branch reuse wins where
many parameters share one schedule, NSR's single combined shift wins where
one parameter touches many non-commuting terms, and the boundary moves
with the term alphabet as the constants predict. Both costs are known
before anything runs: PSR's constant is read off the program, NSR's is
bounded by the static analysis of Sec.~5.3. Each strategy has a
physical imperfection of its own --- PSR the gate's infidelity, NSR the
machine's finite amplitude range --- but neither is a step-size bias: the
gate's enters first order and at the benchmarked fidelity sits well below
FD's floor (Figure~\ref{fig:F6}); the range's appears as rejected draws,
never as bias, unless the series is truncated, where Lemma~D.5 bounds it.
Appendix~G works the regimes case by case.

\begin{figure}[t]
  \centering
  \includegraphics[width=\linewidth]{figs/F_select.pdf}
  \caption{Which sound lowering needs fewer executions, on the design plane
  the compiler reads off program text: $P$ differentiated coefficients
  against $k$, the number of device-alphabet terms each touches. Shading
  marks the winner (blue: PSR; green: NSR); the solid line is the measured
  crossing. Computed at compile time from static certificates alone the
  boundary does not cross this plane---for non-commuting mixed tangents the
  spectral-diameter certificate is loose, so the certificate-guided compiler
  picks PSR everywhere here (the safe direction: looseness costs shots, never
  bias), forfeiting at most $1.9\times$ shots at any sampled point. Star: the
  TFIM instance of Figure~\ref{fig:loop} (per-bond couplings, $k{=}1$); open
  circle: its global-coefficient rewrite. Seven-qubit instance; NSR wins on
  $15\%$ of the sampled plane.}
  \label{fig:select}
\end{figure}

\subsection{At Device Scale}
\label{sec:eval-scale}

\begin{figure*}[t]
  \centering
  \includegraphics[width=\linewidth]{figs/F_scale_combined.pdf}
  \caption{The compiled strategies side by side, and differentiation's price at
  device scale. Left (table): PSR, NSR, and the finite-difference baseline
  compared across three dashed-separated blocks---\emph{Requirements} (top),
  per-gradient \emph{Cost} (middle), and \emph{Guarantees} (bottom). Right (strip): the
  incremental cost of differentiation over the source compilation, log--log, out
  to $10^{3}$ qubits on the target-aware specialized path--- per-branch differentiation is a
  fraction of a percent of the source compile. Differentiation is not the bottleneck.}
  \label{fig:scale}
\end{figure*}

A choice the compiler can only afford at seven qubits is not a choice, so
the last check grows the instance to $10^{3}$ qubits and asks where
compilation spends its time. The answer (Figure~\ref{fig:scale}) is the
source program, and specifically SimuQ's Hamiltonian solver: for arbitrary
Hamiltonians it tops out near a dozen qubits, and for the 
Hamiltonians real machines implement, reaching a thousand took a simple
specialization of the solver. Differentiation itself barely registers on
the same axes, for reasons visible in Sec.~5.3. An NSR branch is the
source schedule with its coefficient arrays shifted: the solver is not
consulted, and the branch is a table written onto an unchanged schedule. A
PSR branch has rigid structure --- the source schedule cut at one time,
with one $\theta$-independent, fixed operation inserted at the
cut --- so the solve is reused on both sides and only the scheduler and
ledger work around the insertion, a fraction of a percent of the source
compile at every size. Since the runtime compiles one branch per execution
(Definition~D.2) and never re-runs the source solve, the bottleneck of a
differentiated program is the bottleneck of the program it differentiates,
and it lies in the part of the stack this paper borrowed rather than
built. Execution cost points the same way --- one inserted operation or
none per branch, where a digital emulation matching the analog route's own
attenuation would pay thousands of gates --- and Sec.~7 takes that up.

The running instance has now been differentiated, compiled, executed under
noise, optimized, re-lowered, and grown by three orders of magnitude, and
at no point did its derivative stop being a program of Sec.~3. That is the
section's claim: the derivative of an analog experiment is a program, and
the stack that compiles the experiment compiles it.

% NOTES: Appendix C.3 = T4 + setpoint model + \nabla C_device. Appendix F =
% device-ready pulse figure(s) at 100+ qubits. Appendix G = F-select
% constants + rho-plane F3. Appendix H = solver breakdown, generic n<=12.
% T2 -> Sec. 7.